\color{unidarkgreen}{\subsection{Автоматическое повторное включение}}
\color{black}

\begin{enumerate}

\item
Автоматическое повторное включение (АПВ) предназначено для быстрого восстановления первоначального состояния электрической сети после аварийного отключения путём автоматического включения выключателя.

\item
Выбор действующих режимов АПВ для выключателя В1 реализуется тремя входными сигналами ФСУ:
\begin{itemize}
    \item \textit{«АПВ Вn ОНТ(АТ) + ННЭn»} (включение с контролем отсутствия напряжения на ошиновке и наличия напряжения на энергообъекте) -- АПВ ошиновки;
    \item \textit{«АПВ Вn ОНЭn + ННТ(АТ)»} (включение с контролем отсутствия напряжения на энергообъекте и наличия напряжения на ошиновке) -- АПВ энергообъекта;
    \item \textit{«АПВ Вn без контролей»} (включение без контроля напряжений).
\end{itemize}

\item
Возможные режимы АПВ приведены в таблице \ref{apv:tbl1}.

\small
\renewcommand{\arraystretch}{1}
\setlength{\extrarowheight}{2pt}

\begin{longtable}{|>{\centering\arraybackslash}m{6.4cm}|>{\centering\arraybackslash}m{3.2cm}|>{\centering\arraybackslash}m{3.2cm}|>{\centering\arraybackslash}m{3.2cm}|}
\caption{Условия включения\hfill\vspace{-0.5\baselineskip}}\label{apv:tbl1}\\
\hline
\rowcolor{gray!30}
& \multicolumn{3}{c|}{Входные сигналы ФСУ} \\
\hhline{>{\arrayrulecolor{gray!30}}->{\arrayrulecolor{black}}|-|-|-|}
\rowcolor{gray!30}
\raisebox{3pt}{\makecell[c]{Режим АПВ}} & 
\makecell[c]{АПВ Вn \\ ОНТ(АТ)+ННЭn} & 
\makecell[c]{АПВ Вn \\ ОНЭn+ННТ(АТ)} & 
\makecell[c]{АПВ Вn \\ без контролей} \\
\hline
\endfirsthead

\caption*{\hspace{3pt}\emph{Продолжение таблицы \ref{apv:tbl1}}\hfill\ }\vspace{-0.5\baselineskip} \\
\hline
\rowcolor{gray!30}
& \multicolumn{3}{c|}{Входные сигналы ФСУ} \\
\hhline{>{\arrayrulecolor{gray!30}}->{\arrayrulecolor{black}}|-|-|-|}
\rowcolor{gray!30}
\raisebox{3pt}{\makecell[c]{Режим АПВ}} & 
\makecell[c]{АПВ Вn \\ ОНТ(АТ)+ННЭn} & 
\makecell[c]{АПВ Вn \\ ОНЭn+ННТ(АТ)} & 
\makecell[c]{АПВ Вn \\ без контролей} \\
\hline
\endhead

\hline
\makecell[c]{КС(УС)} & 0 & 0 & 0 \\
\hline
\makecell[c]{ОНТ(АТ)+ННЭn / КС(УС)} & 1 & 0 & 0 \\
\hline
\makecell[c]{ОНЭn+ННТ(АТ) / КС(УС)} & 0 & 1 & 0 \\
\hline
\makecell[c]{ОНТ(АТ)+ННЭn / ОНЭn+ННТ(АТ) /  КС(УС)} & 1 & 1 & 0 \\
\hline
\makecell[c]{Без контролей} & 0/1 & 0/1 & 1 \\
\hline
\end{longtable}
\normalsize

\item
Как видно из таблицы \ref{apv:tbl1}, режим АПВ с КС(УС) введен всегда и выводится посредством входного сигнала ФСУ \textit{«АПВ Вn без контролей»}. А режимы ОНТ(АТ) и ОНЭn+ННТ(АТ) могут быть введены одновременно.

\item
Пуск АПВ линии осуществляется по цепи несоответствия при аварийном отключении выключателя (по наличию входного сигнала ФСУ \textit{«Сигнал несоотв. Вn»}). Срабатывание происходит при набранной
выдержке времени готовности АПВ с контролем заданного режима и выполнения условий включения. Выдержка времени готовности АПВ после оперативного включения и первого цикла регулируется уставкой
\textbf{«Тгот АПВ Вn»}, после второго цикла -- \textbf{«Твосст гот. АПВ Вn»}.

\item
АПВ может быть выполнено двукратным, ввод второго цикла осуществляется программным ключом \textbf{«2-ой цикл АПВл Вn»}. Первый цикл осуществляется при выполнении условий включения с выдержкой времени \textbf{«Тср АПВ-1 Вn»}. В случае сброса сигнала пуска АПВ при отсчете времени \textbf{«Тср АПВ-1 Вn»} цепь логики сбрасывается, и схема АПВ возвращается в исходное состояние (готовность АПВ с набранной выдержкой готовности). Если первый цикл АПВ был успешным, то начинается набор выдержки времени готовности к повторному действию.

\item
Если после первого цикла АПВ за время \textbf{«Тгот АПВ Вn»} (выдержка времени к повторному действию не успела набраться) повторно поступает сигнал пуска АПВ, то осуществляется 2-ой цикл с выдержкой времени \textbf{«Тср АПВл-2 Вn»}. В случае сброса сигнала пуска АПВ при отсчете времени \textbf{«Тср АПВ-2 Вn»} цепь логики сбрасывается, и схема возвращается в режим набора выдержки времени готовности к повторному действию. Если второй цикл АПВ был успешным, то начинается набор выдержки времени готовности к повторному действию.

\item
Второй цикл АПВ может быть выведен от внешнего сигнала \textit{«Вывод АПВ-2 Вn»}.

\item
Для схем с двумя выключателями на присоединение предусмотрен входной сигнал ФСУ \textit{«Вn ведомый»} для оперативного изменения порядка включения выключателя от АПВ -- \textit{«ведущий»} или \textit{«ведомый»}. По умолчанию выключатель принимается “ведущим”. При наличии сигнала \textit{«Вn ведомый»} время включения от АПВ линии дополнительно увеличивается на выдержку времени \textbf{«Тср АПВ Вn доп.»}.

\item
Длительность импульса включения выключателя от АПВ линии задается уставкой \textbf{«Тимп АПВ Вn»}.

\item
Алгоритм работы АПВ В1 приведен на рисунке \ref{pict_apv:APV}, алгоритм работы АПВ В2 выполнен аналогично. Уставки АПВ В1 приведены в таблице \ref{app_set:apv1}. Уставки АПВ В2 приведены в таблице \ref{app_set:apv2}.

\medskip
\begin{figure}[H]
\centering
\includegraphics[width=0.9\textwidth,height=0.9\textheight,keepaspectratio]{\fbpath/05.02 ДЗО/761.3041 АПВ/img/APV.pdf}
\captionof{figure}{Функциональная схема логики АПВ В1}\label{pict_apv:APV}
\end{figure}



\end{enumerate}